\documentclass{umbruch-position}

\positionsnummer{ST-001 ES}
\title{Executive Summary}
\untertitel{Systemtheoretische Aufarbeitung der Regressangst --- ST-001 v0.15}
\author{Um:bruch -- Denkfabrik f\"ur gesellschaftlichen Wandel}
\date{April 2026}
\versionsnummer{1.2}

\zusammenfassung{%
Das Wirtschaftlichkeitspr\"ufsystem nach \S\S\,106\,ff. SGB\,V erzeugt eine Regressangst, deren Folgekosten die durchgesetzten Regresse nach Modellrechnung um das 100- bis 1.000-fache \"ubersteigen k\"onnten. Die Studie ST-001 identifiziert 15 Befunde, die diese Diskrepanz aus historischen, rechtsdogmatischen, verhaltens\"okonomischen und international vergleichenden Quellen aufkl\"aren. \textbf{Genre und Geltungsanspruch:} ST-001 ist keine Original-Studie im Sinne einer neuen empirischen Prim\"arerhebung, sondern eine \textbf{multiperspektivische Bestandsaufnahme mit Pilot-Charakter} --- eine Diagnose-Konstruktion (drei Zertifikatsfehler) auf Basis integrierter rechtsdogmatischer, statistischer, verhaltens\"okonomischer und institutionentheoretischer Perspektiven. 15 Befunde, drei minimal-invasive Reformhebel und 15 bewertete Ma\ss{}nahmen (KNV-Tabelle) bieten konkrete Handlungsoptionen f\"ur Politik, Verb\"ande und Zivilgesellschaft.%
}

\schlagworte{Regress, Wirtschaftlichkeitspr\"ufung, V2b, Transparenz, Executive Summary}

\begin{document}
\emergencystretch=3em
\maketitle

%% ============================================================
\section{Forschungsfrage}
%% ============================================================

Die Studie untersucht eine auff\"allige Diskrepanz im deutschen Wirtschaftlichkeitspr\"ufsystem nach \S\S\,106\,ff. SGB\,V: Die offizielle Regressdurchsetzungsquote liegt unter 1\,\% der Vertrags\"arzte --- dennoch berichten 72\,\% der Haus\"arzte, mindestens einmal einen Regress erlebt zu haben (Ribbat et al. 2023, n=770), und 47\,\% geben an, die Regressgefahr besch\"aftige ihren Praxisalltag \glqq{}stark oder sehr stark\grqq{}. Die zentrale Frage lautet: \textbf{Was passiert im System, das beides zugleich produziert?} Die Studie versucht, diese Diskrepanz aus historischen, rechtsdogmatischen, verhaltens\"okonomischen, gesundheits\"okonomischen und international vergleichenden Quellen aufzukl\"aren.

ST-001 ist die wissenschaftliche Begleituntersuchung zum Konzeptpapier PP-003 (Regress-Transparenzportal).

%% ============================================================
\section{Zentrale Befunde}
%% ============================================================

Die Studie identifiziert 15 wesentliche Befunde, die hier mit ihrer jeweiligen Evidenzbasis und ihren Einschr\"ankungen dargestellt werden.

% --------------------------------------------------
\subsection{Vier Messebenen und ihre Vermischung}

\textbf{Befund:} Der Begriff \glqq{}Regressquote\grqq{} wird in der \"offentlichen Diskussion f\"ur vier grundlegend verschiedene Messgr\"o\ss{}en verwendet: Pr\"ufquote, Regressverfahrenquote, Regressdurchsetzungsquote und Regressvolumen. Diese vier Ebenen sind nicht miteinander vergleichbar, werden aber regelm\"a\ss{}ig vermischt. Die \"offentlich kommunizierte \glqq{}unter 1\,\%\grqq{}-Zahl bezieht sich auf die Regressdurchsetzungsquote auf Arztebene, \"uber beide Verfahren (V1+V2) hinweg, ohne Verfahrenstrennung. Die 99,7-\%-Verwerfungsrate (KV Hessen 2015) misst hingegen die Pipeline-Verwerfungsrate der Auff\"alligkeitspr\"ufung (V1) allein.

\textbf{Evidenzbasis:} KV-Statistiken (KV Hessen 2015, Dt. \"Arzteblatt Art. 113630; KV BaW\"u 2022), eigene systematische Begriffsanalyse.

\textbf{Einschr\"ankung:} Bundesweit fehlen alle vier Ebenen als einheitlich ver\"offentlichte Zahlen. Keine KV publiziert verfahrensgetrennte (V1 vs. V2) Statistiken.

% --------------------------------------------------
\subsection{Definitionsdivergenz: Warum ``unter 1\,\%'' und ``72\,\%'' kein Widerspruch sind (CORE4)}

\textbf{Befund:} Der Begriff ``Regress'' wird von verschiedenen Akteuren unterschiedlich definiert. Diese Definitionsdivergenz erkl\"art die scheinbare Diskrepanz:

\smallskip
{\small
\begin{tabularx}{\textwidth}{l p{4.5cm} l l}
\hline
\textbf{Schicht} & \textbf{Was gez\"ahlt wird} & \textbf{Zahl} & \textbf{Quelle} \\
\hline
1 (eng) & Nur V1, rechtskr\"aftig & ${<}$100/J $\to$ 0,065\,\% & Virchowbund 2018 \\
2 (amtlich) & V1+V2, rechtskr\"aftig & 21/21.500 $\to$ 0,097\,\% & KV BaW\"u 2022 \\
3 (faktisch) & Alle V2 inkl. Vergleiche & ca.\,47.000/J $\to$ ${\sim}$30\,\% & BMG GVSG 2024 \\
\hline
\end{tabularx}
}
\smallskip

Der Faktor zwischen Schicht~1 und Schicht~3 betr\"agt ca.\,470--500$\times$. Alle Akteure sagen die Wahrheit --- \"uber verschiedene Messgr\"o\ss{}en. Die Diskrepanz ist \textbf{kein Wahrnehmungsfehler der \"Arzte, sondern ein Messproblem der Statistik}. In Bayern meldet der SpiFa 13.332 festgesetzte Regresse (2024) bei $\sim$24.000 Vertrags\"arzten --- Durchschnittsbetrag 334\,EUR, $\sim$70\,\% davon Bagatelle unter 300\,EUR.

\textbf{Evidenzbasis:} BMG GVSG-Referentenentwurf 2024; SpiFa Bayern 2023/2024; Virchowbund 2018; KV BaW\"u 2022; Multi-Agenten-Recherche (15 Agenten, 25 Datenpunkte, 14 Domains).

\textbf{Einschr\"ankung:} Die Sch\"atzung ``${\sim}$30\,\%'' basiert auf der Annahme gleichm\"a\ss{}iger Verteilung der 47.000 Verfahren auf 155.000 \"Arzte. Die tats\"achliche Verteilung (nach Fachgruppe, Region, Kassenart) ist nicht publiziert.

\textbf{Originalit\"at:} Nach systematischem Prior-Art-Check (Opus-Agent, 20+ Suchanfragen) ist die Synthese --- Definitionsdivergenz als Erkl\"arung f\"ur die Ribbat-Diskrepanz --- \textbf{nirgends publiziert}. Die gesamte Literatur erkl\"art die Diskrepanz als psychologisches Problem der \"Arzte (Availability Bias, irrationale Angst). Alle Einzelbausteine (Ribbat, BMG, SpiFa, Virchowbund) sind publiziert; die Verbindung ist origin\"ar.

% --------------------------------------------------
\subsection{Zwei Pr\"ufdom\"anen (V1/V2) und ihre unterschiedliche Fairness}

\textbf{Befund:} Das Wirtschaftlichkeitspr\"ufsystem besteht aus zwei funktional inkommensurablen Verfahren:

\begin{itemize}
  \item \textbf{V1 (Auff\"alligkeitspr\"ufung, \S\,106 SGB\,V):} Statistisches Screening, initiiert von der Pr\"ufungsstelle, ohne individuellen Verdacht. Breiter Filter mit hoher Sensitivit\"at und niedriger Spezifit\"at --- eine hohe Verwerfungsrate (99,7\,\%) ist die normale Funktionsweise.
  \item \textbf{V2 (Einzelfallpr\"ufung, \S\,106b SGB\,V):} Antragsgebundenes Verdachtsverfahren, initiiert von der Krankenkasse. Gesetzlich ohne \glqq{}Beratung vor Regress\grqq{} (\S\,106b Abs.\,2 S.\,4 SGB\,V). Kasse ist gleichzeitig Antragstellerin und Empf\"angerin des Regressbetrags.
\end{itemize}

V1 hat \"uber die Jahre Schutzschichten dazubekommen (Beratung vor Regress 2012, Praxisbesonderheiten 2011, Richtgr\"o\ss{}en abgeschafft 2015). V2 blieb explizit ohne Beratungsschutz --- das war kein \"Ubersehen, sondern bewusstes gesetzliches Design.

\textbf{Evidenzbasis:} Gesetzestexte (SGB\,V, GKV-VSG 2015), BSG-Rechtsprechung, KBV-Stellungnahme zum GKV-VSG 2015, regionale Pr\"ufvereinbarungen.

\textbf{Einschr\"ankung:} V2-Daten existieren (BMG: ca.\,47.000 Verfahren bundesweit 2022; SpiFa Bayern: 13.000 Regresse/Jahr), werden aber nicht als integrierte Durchsetzungsstatistik mit Verwerfungs- und Vergleichsquoten publiziert. Keine KV publiziert verfahrensgetrennte (V1 vs. V2) Pipeline-Daten.

% --------------------------------------------------
\subsection{V2a vs. V2b --- die entscheidende Differenzierung (CORE3), V2b-Binnendifferenzierung (v0.14)}

\textbf{Befund:} Die Einzelfallpr\"ufung (V2) zerf\"allt in zwei Unterkategorien mit fundamental verschiedener Verfahrenslogik, V2b ist zus\"atzlich binnendifferenziert:

\begin{itemize}
  \item \textbf{V2a (Unwirtschaftliche Verordnung):} Leitlinienkonform, aber teurer als n\"otig. Medizinische Verteidigung m\"oglich; die Kostendifferenzmethode (\S\,106b Abs.\,2a SGB\,V, BSG 2024) begrenzt den Regress auf die Differenz zur wirtschaftlichen Alternative.
  \item \textbf{V2b (Unzul\"assige Verordnung):} Medizinische Gr\"unde grunds\"atzlich irrelevant (normativer Schadensbegriff, BSG B\,6\,KA\,5/09\,R); keine Kostendifferenzmethode; keine Heilbarkeit nach Bezahlung durch die Kasse; kein Beratungsschutz. Intern binnendifferenziert:
  \begin{itemize}
    \item \textbf{V2b\,(i) --- reine Formfehler} (fehlende Unterschrift, Stempel statt Unterschrift, fehlende ICD, falsches Aut-idem-Kreuz): Formkonformit\"at ist kein Proxy f\"ur medizinische Qualit\"at. Voller Zertifikatsfehler.
    \item \textbf{V2b\,(ii) --- inhaltliche Unzul\"assigkeit} (AM-RL-Versto\ss{}, Off-Label ohne Begr\"undung, AMNOG-Ausschluss): Die Formregel codiert bereits medizinische G-BA-Evidenz; der Zertifikatsfehler ist schw\"acher, weil die Regel medizinische Qualit\"at \emph{mittelbar} misst --- das Problem ist die punitive Kommunikation statt eines Lernsignals.
  \end{itemize}
\end{itemize}

\textbf{Konsequenz f\"ur das Gesamtbild:}
\begin{itemize}
  \item V1 ist \textbf{fairer als in fr\"uhen Versionen der Studie dargestellt} (Anh\"orung, medizinische Ber\"ucksichtigung, Stufenmodell).
  \item V2a hat \textbf{Verteidigungsm\"oglichkeiten} (medizinische Begr\"undung, Kostendifferenz).
  \item V2b ist \textbf{isoliert betrachtet noch problematischer} als bisher dargestellt: Es ist das einzige Verfahren im deutschen Recht ohne Heilungsm\"oglichkeit, ohne Kostendifferenz, ohne medizinische Relevanz.
\end{itemize}

\textbf{Quantitative Einsch\"atzung (v1.2).} Nicht \emph{alle} V2-Verfahren sind dysfunktional. Gesch\"atzt \textbf{60--65\,\% der V2-Verfahren} (Kategorie~B reine Formfehler + verfahrensfremde Angriffsvektoren wie der Stempel-Regress) fallen unter den zweiten Zertifikatsfehler; die restlichen 35--40\,\% (Kategorie~A: evidenzbasierte Formregeln --- AM-RL-Anlagen, G-BA-Beschl\"usse, Biosimilar-Pflicht) codieren medizinische Evidenz --- allerdings \"uber den punitiven statt den edukativen Kanal. Die Forderung der Studie ist deshalb nicht die Abschaffung der Evidenzsteuerung, sondern die Erg\"anzung des punitiven Kanals um einen edukativen (Beratungsschutz auch bei V2, Vorabinformation, Hebel~1).

\textbf{Evidenzbasis:} GPE BaW\"u, KV Nordrhein, KV Bayern, BSG-Rechtsprechung (B\,6\,KA\,5/09\,R, B\,6\,KA\,5/23\,R, B\,6\,KA\,27/12\,R), eigene Detailrecherche (CORE3-Dokument).

\textbf{Einschr\"ankung:} Die V2a/V2b-Differenzierung wird in keiner offiziellen Statistik abgebildet. Welcher Anteil der V2-Verfahren V2a und welcher V2b betrifft, ist nicht \"offentlich bekannt.

% --------------------------------------------------
\subsection{Drei Zertifikatsfehler --- doppelt im Messinhalt, einmal in der Messrichtung}

\textbf{Befund:} Das Regresssystem weist drei Zertifikatsfehler auf --- zwei im Messinhalt, einen \"ubergeordneten in der Messrichtung:

\begin{enumerate}
  \item \textbf{Erster Zertifikatsfehler --- V1 (Kostenabweichung $\neq$ Medizin):} Die Auff\"alligkeitspr\"ufung misst Kostenabweichung vom Fachgruppendurchschnitt, nicht medizinische Qualit\"at. F\"ur ein Screening ist das technisch korrekt; der Fehler liegt in der \"offentlichen Interpretation der niedrigen Endquote als Beleg f\"ur ein \glqq{}funktionierendes System\grqq{}. Das Schutzstufenmodell filtert fast vollst\"andig.
  \item \textbf{Zweiter Zertifikatsfehler --- V2b (Form $\neq$ Medizin):} Die Einzelfallpr\"ufung beansprucht, medizinische Unwirtschaftlichkeit zu pr\"ufen. Sie misst faktisch formale bzw.\ zulassungsrechtliche Angreifbarkeit. Der Befund gilt in voller Sch\"arfe f\"ur V2b\,(i) (reine Formfehler); bei V2b\,(ii) (inhaltliche Unzul\"assigkeit, AM-RL-Versto\ss{}) ist er abgeschw\"acht, weil die Formregel G-BA-Evidenz codiert --- das Problem bleibt die punitive statt edukative Kommunikation.
  \item \textbf{\"Ubergeordneter Zertifikatsfehler --- Einbahnpr\"ufung (Richtung $\neq$ Symmetrie):} Das gesamte System pr\"uft ausschlie\ss{}lich \glqq{}zu viel verordnet\grqq{}, nie \glqq{}zu wenig versorgt\grqq{}. Der Versorgungsauftrag des \S\,70 SGB\,V hat kein eigenes Messinstrument im ambulanten Bereich. Damit verschiebt das System den Schaden vom Arzt zum Patienten: Der Arzt ist bei Handeln dem Regressrisiko ausgesetzt; der Patient tr\"agt bei \"arztlicher Zur\"uckhaltung die Unterversorgung und bei \"arztlicher Ausweichreaktion die \"Uberversorgung (iatrogene Risiken + finanziell/zeitliche Zuzahlungen, Fahrtwege, Arbeits-/Urlaubszeit). SVR-Trias: Unter-/\"Uber-/Fehlversorgung (SVR 2000/2001, \S\,135a SGB\,V).
\end{enumerate}

Medizinische Angemessenheit im Sinne des \S\,12 Abs.\,1 SGB\,V und der Versorgungsauftrag des \S\,70 SGB\,V werden systematisch nur au\ss{}erhalb des Regresssystems zertifiziert (Schlichtungsstellen der \"Arztekammern, zivilrechtliche Arzthaftung, keine operative Versorgungspr\"ufung im ambulanten Bereich).

\textbf{Evidenzbasis:} Systematische Analyse aller Pr\"ufstufen (Pipeline-Diagramme); BSG-Rechtsprechung; \S\S\,12, 70, 135a SGB\,V; Goodhart's Law / Campbell's Law; SVR-Gutachten 2000/2001.

\textbf{Einschr\"ankung:} Der Einwand, das System habe nie ein medizinisches Qualit\"atsversprechen gegeben, ist dogmatisch zutreffend (Einwand 2 der Studie). Der Zertifikatsfehler beschreibt nicht einen Fehler \textit{im} System, sondern die L\"ucke zwischen Systemmessung und \"offentlicher Bedeutungszuschreibung. V1 ber\"ucksichtigt medizinische Besonderheiten im Schutzstufenmodell; V2a l\"asst medizinische Verteidigung zu; V2b\,(ii) codiert G-BA-Evidenz. Der \"ubergeordnete Richtungsfehler gilt f\"ur \emph{alle} Verfahren gleicherma\ss{}en.

% --------------------------------------------------
\subsection{V1-Schutzstufenmodell}

\textbf{Befund:} Die Auff\"alligkeitspr\"ufung (V1) enth\"alt ein mehrstufiges Schutzmodell, das in der \"offentlichen Debatte bisher wenig gew\"urdigt wurde:

\smallskip
{\small
\begin{tabularx}{\textwidth}{lX}
\hline
\textbf{Stufe} & \textbf{Schutzwirkung} \\
\hline
Anh\"orung / Stellungnahme & Arzt kann Abweichung erkl\"aren (6 Wochen Frist) \\
Praxisbesonderheiten & Medizinische Besonderheiten der Patientenstruktur werden ber\"ucksichtigt \\
Beratung vor Regress & Bei erstmaliger Auff\"alligkeit: Beratung statt Regress (\S\,106b Abs.\,2) \\
Karenzzeit & Erst im n\"achsten Zeitraum pr\"ufbar \\
Deckel & 25.000\,EUR in den ersten zwei Jahren nach Beratung \\
5-Jahres-Regel & Nach 5 Jahren ohne Auff\"alligkeit gilt man wieder als erstmalig \\
\hline
\end{tabularx}
}
\smallskip

Die 99,7-\%-Verwerfungsrate ist das Ergebnis eines \textbf{funktionierenden Schutzstufenmodells}, nicht eines dysfunktionalen Filters.

\textbf{Evidenzbasis:} GPE BaW\"u, KV Nordrhein, \S\,106b Abs.\,2 SGB\,V.

\textbf{Einschr\"ankung:} Alle markierten \"Arzte tragen dennoch die psychologischen Kosten des Screenings (Angst, Aufwand), auch die 99,7\,\%, die entlastet werden. Praxisbesonderheiten m\"ussen aktiv geltend gemacht werden --- wer das nicht wei\ss{}, verliert.

% --------------------------------------------------
\subsection{Ribbat 72\,\% --- Regresserfahrung ohne Verfahrenstrennung}

\textbf{Befund:} 72\,\% der Haus\"arzte und 78\,\% der Orthop\"aden berichten, mindestens einmal einen Regress erlebt zu haben. 47\,\% (HA) / 55\,\% (Ortho) geben an, die Regressgefahr besch\"aftige den Praxisalltag \glqq{}stark oder sehr stark\grqq{}. 51\,\% der Haus\"arzte \"uberweisen \glqq{}zumindest gelegentlich trotz Indikation\grqq{} an Spezialisten. 85\,\% haben mindestens einmal aus Regressangst eine Verordnung unterlassen.

\textbf{Evidenzbasis:} Ribbat L, Linde K, Schneider A, Riedl B (2023): Gesundheitswesen 85(2): 111--118. DOI: 10.1055/a-1594-2527. Bundesweite Zufallsstichprobe, je 1.000 Haus\"arzte und 1.000 Orthop\"aden, R\"ucklauf 41\,\% / 39\,\% (n=770). Lehrstuhl f\"ur Allgemeinmedizin TU M\"unchen.

\textbf{Einschr\"ankung:} Die Studie misst \textbf{V1+V2 zusammen, ohne Verfahrenstrennung}. Sie misst subjektive Wahrnehmung und selbstberichtetes Verhalten, keine objektiven Versorgungsergebnisse. Eine unabh\"angige Replikation steht aus. Der kausale Schluss \glqq{}Pr\"ufdruck $\to$ Unterversorgung\grqq{} ist nicht empirisch geschlossen (Einwand 1 der Studie). Die Studie ist der beste verf\"ugbare Datenpunkt, nicht der abschlie\ss{}ende Beweis.

\textbf{Rechnerische Plausibilisierung:} Ein Poisson-Modell mit drei Ankern ($\lambda = 0{,}30$ aus der BMG-Bundesrechnung, $\lambda = 0{,}85$ aus der KV-Nordrhein-Zahlungsquote, $\lambda = 1{,}48$ aus der KV-Nordrhein-Antragsquote) ergibt eine Bandbreite: Eine deutliche Mehrheit der \"Arzte (60--99\,\% \"uber drei Jahre, je nach gez\"ahltem Ereignistyp --- abgeschlossenes Verfahren, Zahlung oder Antrag) ist betroffen. Die 72\,\% Ribbat-Erfahrungsrate ist damit rechnerisch plausibel, unabh\"angig von der ``unter 1\,\%''-Durchsetzungsquote (vgl.\ ST-001 Abschn.~3.4, Summierungsrechnung).

% --------------------------------------------------
\subsection{Spieltheoretisches Nash-Gleichgewicht}

\textbf{Befund:} Bei einem wahrgenommenen Regressrisiko von ca. 15\,\% ($P_{\text{perceived}}$) wird \glqq{}nicht verordnen\grqq{} zur dominanten Strategie des Arztes. Das Ergebnis ist ein Pareto-suboptimales Nash-Gleichgewicht: Arzt verliert (ethischer Stress, Mehrarbeit), Patient verliert (Unterversorgung), Kasse gewinnt nichts (zahlt sp\"ater die Hospitalisierung), System verliert die Folgekosten. Niemand will dieses Gleichgewicht --- alle spielen es individuell rational.

Zus\"atzlich identifiziert die Studie ein Solidarit\"ats-Gefangenendilemma innerhalb der \"Arzteschaft: Wer \"offentlich k\"ampft, wird zur Zielscheibe. Wer schweigt, profitiert vom Erfolg anderer.

\textbf{Evidenzbasis:} Verhaltens\"okonomische Modellierung auf Basis der Ribbat-Daten, Kahneman/Tversky-Verlustaversion ($\lambda \approx 2{,}25$), KBV-Berufsmonitoring.

\textbf{Einschr\"ankung:} Die Schwellenwerte sind Modellannahmen. Die formale Durchmodellierung mit Auszahlungsmatrizen und Sensitivit\"atsanalysen ist als Folgeanalyse vorgesehen, liegt aber noch nicht vor.

% --------------------------------------------------
\subsection{Folgekostenrechnung (0,9--1,7 Mrd. EUR, Szenario-Annahme)}

\textbf{Befund:} Die Studie sch\"atzt die Folgekosten regressangst-induzierter Fehlversorgung im Dreistufenmodell:

\smallskip
{\small
\begin{tabularx}{\textwidth}{lXl}
\hline
\textbf{Stufe} & \textbf{Wert} & \textbf{Quelle} \\
\hline
Systemrahmen (ACSC + Non-Adh\"arenz, vor Attribution) & 7--17 Mrd. EUR/Jahr & WIdO, Zi, DIMDI HTA 65 \\
$\times$ Attributionsrate (konservativ) & 5\,\% & Szenario-Annahme \\
= Regress-Folgekosten (Gesamtbandbreite) & \textbf{0,9--1,7 Mrd. EUR/Jahr} & Eigene Hochrechnung \\
davon mittleres Szenario & \textbf{rund 1,1 Mrd. EUR/Jahr} & Sensitivit\"atstabelle ST-001 Abschn.~4.3 \\
\hline
\end{tabularx}
}
\smallskip

Dem stehen durchgesetzte Regresse von gesch\"atzt 5--50 Mio. EUR/Jahr gegen\"uber. \textbf{Szenario-Verh\"altnis Schaden zu Einnahme: ca. 100:1 bis 1.000:1.}

Im mittleren Szenario w\"aren ca. 6.230 vermeidbare Hospitalisierungen und ca. 125 Todesf\"alle pro Jahr anteilig mit regressangst-induzierter Unterversorgung assoziierbar.

\textbf{Evidenzbasis:} WIdO Krankenhaus-Report, Zi 2021, DIMDI HTA 65, OECD Health at a Glance 2025, Ribbat 2023, eigene Hochrechnung mit Sensitivit\"atstabelle.

\textbf{Einschr\"ankung:} \textbf{Diese Zahlen sind ein illustratives Gedankenexperiment auf Basis von Szenario-Annahmen, kein kausalanalytisch abgesicherter Befund.} Insbesondere die Attributionsrate (5\,\%) und die Frequenz unterlassener Verordnungen sind Modellannahmen. Keine der verf\"ugbaren Studien misst diese Gr\"o\ss{}en direkt. Der unvollst\"andige Nenner: Die Abschreckungswirkung des Systems (verhinderte Mehrausgaben) ist real, aber nicht quantifiziert. Die 100:1-Zahl bildet den worst case f\"ur das System ab, nicht den Erwartungswert.

% --------------------------------------------------
\subsection{Normativer Schadensbegriff und V2b-Unheilbarkeit}

\textbf{Befund:} Der normative Schadensbegriff des BSG (B\,6\,KA\,5/09\,R, 2010) definiert: Der Schaden der Kasse gilt als ab Bezahlung der formal fehlerhaften Verordnung eingetreten --- unabh\"angig davon, ob das Medikament medizinisch indiziert war. Nach Verfahrenseinleitung ist der Formfehler nicht mehr heilbar. Das BSG hat dies 2025 (B\,6\,KA\,9/24\,R) bei 491.000\,EUR Stempel-Regress best\"atigt: Das Argument der medizinischen Indikation wurde als irrelevant erkl\"art.

2024 schloss das BSG die Differenzkostenregelung f\"ur \glqq{}unzul\"assige\grqq{} Verordnungen (= Formfehler) ausdr\"ucklich aus (B\,6\,KA\,5/23\,R; B\,6\,KA\,10/23\,R).

\textbf{Evidenzbasis:} BSG-Rechtsprechung (direkt verifiziert).

\textbf{Einschr\"ankung:} In seltenen F\"allen haben Gerichte trotz Formfehler zugunsten des Arztes entschieden (SG Mainz 2022: 260.000\,EUR als rechtsmissbr\"auchlich zur\"uckgewiesen; LSG BaW\"u 2025: Doppelpr\"ufungsverbot). Diese Ausnahmen belegen durch ihre Seltenheit, wie wenig Schutz der Rechtsweg typischerweise bietet.

% --------------------------------------------------
\subsection{Fehlversorgung bidirektional (Unter- UND \"Uberversorgung) --- SVR-Trias}

\textbf{Befund:} Die deutsche Regressangst erzeugt --- anders als das US-amerikanische Tort Law --- eine Fehlversorgung in zwei Richtungen zugleich (Begrifflichkeit: SVR 2000/2001, \S\,135a SGB\,V):

\begin{itemize}
  \item \textbf{Unterversorgung bei Verordnungen} (\S\,70 SGB\,V): \"Arzte verschreiben weniger, vermeiden teure Medikamente, verzichten auf Off-Label-Therapien.
  \item \textbf{\"Uberversorgung bei Diagnostik und \"Uberweisungen} (\S\,12 SGB\,V): \"Arzte veranlassen zus\"atzliche Bluttests, \"uberweisen an Fach\"arzte, fordern Bildgebung an --- nicht weil der Patient es braucht, sondern zur eigenen Absicherung. (Goetz 2024: 54\,\% w\"ochentlich unn\"otige Labortests; Ribbat 2023: 51\,\% \"uberweisen trotz eigener Kompetenz.)
\end{itemize}

Der pr\"azisere Begriff ist \textbf{Fehlversorgung} als Querkategorie: Das System lenkt \"arztliche Ressourcen von der Therapie in die Absicherung. \textbf{Schadens-Verschiebung zum Patienten} --- in beide Richtungen: (a) \textit{medizinisch} iatrogene Risiken und Overdiagnosis-Belastung; (b) \textit{finanziell und zeitlich} zus\"atzliche Arzttermine, Zuzahlungen, Wartezeiten, Urlaubstage, Anfahrtswege, Begleitpersonen. Was das Prüfsystem als ``Wirtschaftlichkeit'' einspart, zahlt der Patient doppelt --- als verz\"ogerte Therapie \emph{und} als Zeit-/Wegelast unn\"otiger Zusatzleistungen.

\textbf{Evidenzbasis:} Ribbat 2023, Goetz 2024, internationaler Vergleich (Kessler \& McClellan 1996, Currie \& MacLeod 2008, Mello et al. 2010).

\textbf{Einschr\"ankung:} Unterversorgung ist methodisch viel schwerer zu messen als \"Uberversorgung. Wenn ein Arzt eine notwendige Verordnung unterl\"asst, taucht es nirgends auf --- au\ss{}er sp\"ater in der Hospitalisierungsstatistik, ohne Ursachenzuschreibung.

% --------------------------------------------------
\subsection{V2-Durchsetzung im Vergleich versteckt (70--80\,\%)}

\textbf{Befund:} Gesch\"atzt 70--80\,\% der von V2 betroffenen \"Arzte schlie\ss{}en Vergleiche (typisch 30--50\,\% des Regressbetrags). Der Vergleich ist die Durchsetzung --- nur au\ss{}erhalb des Gerichts. Er erscheint in der Statistik als \glqq{}kein Regress\grqq{}, obwohl der Arzt gezahlt hat. Die Verwerfungsquote misst damit nicht die Qualit\"at der Antr\"age, sondern die Gegenwehrst\"arke der Betroffenen.

Nur zwei von f\"unf Beendigungswegen produzieren eine inhaltliche Feststellung (Ablehnungsbescheid oder Gerichtsurteil). Da fast alle Verfahren vor Gericht enden, bleibt die Systemqualit\"at dauerhaft ungemessen.

\textbf{Evidenzbasis:} Praxisberichte von Regressberatern und \"arztlichen Fachverb\"anden, Verfahrensstrukturanalyse.

\textbf{Einschr\"ankung:} Es existiert \textbf{keine publizierte Studie} zur Vergleichsquote bei V2. Die 70--80-\%-Sch\"atzung beruht auf Praxisberichten, nicht auf systematischen Erhebungen.

% --------------------------------------------------
\subsection{Container-Genealogie (9 Normcontainer)}

\textbf{Befund:} 1989 reichte ein einziger Paragraph (\S\,106 SGB\,V) f\"ur die gesamte Wirtschaftlichkeitspr\"ufung. 2026 sind daraus f\"unf ambulante (\S\S\,106, 106a, 106b, 106c, 106d), zwei station\"are (\S\S\,275, 275c), \S\,48 BMV-\"A plus Rahmenvorgaben und bis zu 17 regionale Pr\"ufvereinbarungen geworden. Die geschriebene Regelfl\"ache hat sich in 35 Jahren etwa verachtfacht. Der Regelungsgehalt ist im Kern derselbe geblieben.

Drei methodisch besonders aufschlussreiche Wanderungen: (1) Die \S\,106a-Umbenennungsfalle (2017), (2) die Auslagerungskaskade der Einzelfallpr\"ufung auf drei Ebenen, (3) die parallel-gegenl\"aufige Operation ambulant/station\"ar. Sieben von zehn identifizierten strukturellen Inkonsistenzen im Normgef\"uge w\"aren mit minimalem gesetzgeberischem Aufwand fixbar. Sie werden seit Jahren nicht behoben.

\textbf{Evidenzbasis:} Gesetzestexte (GRG 1989, GSG 1992, GKV-VSG 2015, TSVG 2019), BT-Drucksachen, historische Fassungen via BGBl.

\textbf{Einschr\"ankung:} Die Containeranalyse zeigt Komplexit\"atszunahme, nicht kausal Dysfunktion. Die praktischen Folgen lassen sich nur \"uber die Befunde aus Teil III (Verhaltens\"okonomie) und Teil V (institutionelles Schweigen) der Studie ableiten.

% --------------------------------------------------
\subsection{Internationaler Vergleich}

\textbf{Befund:} Deutschland und die Schweiz sind in Westeuropa praktisch die einzigen L\"ander mit fl\"achendeckender individueller Regress-Pr\"ufung gegen einzelne \"Arzte. Frankreich entwickelt seit 2023 eine strukturell parallele Debatte unter dem Begriff \glqq{}d\'elit statistique\grqq{} (MSO/MSAP-Verfahren der CNAM). Andere L\"ander (UK, NL, Skandinavien) steuern \"arztliches Verordnungsverhalten \"uber positive Anreize (QOF), Capitation, zentrale Preisbildung oder professionelle Selbstregulation --- ohne Individual-Sanktionen.

\textbf{Evidenzbasis:} OECD Health at a Glance, NHS-QOF-Dokumentation, franz\"osische Prim\"arquellen (MG France, FMF, Le Quotidien du M\'edecin, CNAM-Kontrollzahlen 2024).

\textbf{Einschr\"ankung:} Die Vergleichbarkeit ist durch unterschiedliche Systemarchitekturen begrenzt. Der franz\"osische Fall betrifft prim\"ar Krankschreibungen, nicht Arzneimittel.

% --------------------------------------------------
\subsection{Goodhart's Law / Campbell's Law}

\textbf{Befund:} Das Proxy-Versagen des Pr\"ufsystems --- Formkonformit\"at wird gemessen statt medizinischer Qualit\"at --- ist kein Einzelfall, sondern eine bekannte systemische Pathologie. Goodhart's Law (\glqq{}When a measure becomes a target, it ceases to be a good measure\grqq{}) und Campbell's Law beschreiben exakt diesen Mechanismus. Die Pr\"ufquote wird zum Ziel, die Formkonformit\"at zum Messobjekt --- und beide verlieren dadurch ihre diagnostische Aussagekraft.

\textbf{Evidenzbasis:} Theoriegest\"utzte Analyse, angewandt auf die empirischen Befunde aus Teilen I--III der Studie.

\textbf{Einschr\"ankung:} Die Anwendung ist analytisch, nicht empirisch getestet.

% --------------------------------------------------
\subsection{Doppelfunktion: Evidenz in formaler Verkleidung (CORE4)}

\textbf{Befund:} Die bisherige Analyse behandelte ``Formalit\"at'' als einheitlich medizinisch irrelevant. Das war zu pauschal. Die V2-Pr\"ufanl\"asse zerfallen in drei Kategorien:

\smallskip
{\small
\begin{tabularx}{\textwidth}{l l X}
\hline
\textbf{Kat.} & \textbf{Anteil} & \textbf{Charakter} \\
\hline
A & ${\sim}$35\,\% & Rationale Evidenzsteuerung (Biosimilar-Pflicht, Off-Label, AM-RL) \\
B & ${\sim}$40\,\% & Reine Formalkontrolle (Stempel, ICD-Kodierung) \\
C & ${\sim}$25\,\% & B\"urokratie-Overhead (Bagatelle ${<}$300\,EUR) \\
\hline
\end{tabularx}
}
\smallskip

Das System hat zwei parallele Evidenzkan\"ale: Fortbildungspflicht (CME, \S\,95d SGB\,V, \textit{pr\"aventiv}) und Einzelfallpr\"ufung (V2, \textit{punitiv}). Der Arzt erh\"alt kein Lernsignal vor dem Regress. Die Forderung verschiebt sich: Nicht Abschaffung der Evidenzsteuerung (Kat.\,A), sondern Erg\"anzung des punitiven Kanals um einen edukativen.

\textbf{Evidenzbasis:} AOK BaW\"u Pr\"ufthemenliste 2025; SpiFa Bayern; AM-RL-Struktur; \S\,95d SGB\,V.

\textbf{Einschr\"ankung:} Die Prozentanteile (35/40/25) sind Sch\"atzungen auf Basis der Pr\"ufthemenlisten, nicht empirisch gemessen. Eine verfahrensgetrennte Aufschl\"usselung nach V2a/V2b existiert nicht.

\textbf{Originalit\"at:} Prior-Art-Check (Opus-Agent, 20+ Suchanfragen): Die Synthese --- Regress + CME als zwei parallele Evidenzkan\"ale (punitiv vs. edukativ) --- ist nirgends publiziert. Stallberg, IQWiG, Zeschick und Gandjour thematisieren Einzelaspekte; die Verbindung ist origin\"ar.

% --------------------------------------------------
\subsection{Foucault-Panoptikum}

\textbf{Befund:} Das Wirtschaftlichkeitspr\"ufsystem erf\"ullt die drei Kernmerkmale des Foucaultschen Panoptikums: (1)~Asymmetrische Sichtbarkeit --- der Arzt wei\ss{} nicht, ob und wann seine Verordnungen gepr\"uft werden, die Kasse sieht seine Daten in Echtzeit. (2)~Internalisierung der Kontrolle --- 72\,\% der Haus\"arzte \"andern ihr Verordnungsverhalten pr\"aventiv, nicht weil sie gepr\"uft werden, sondern weil sie gepr\"uft werden \textit{k\"onnten}. (3)~\"Okonomie der Macht --- die niedrige V1-Regressquote ist maximale Effizienz: Das Screening allein erzeugt fl\"achendeckende Verhaltens\"anderung, ohne viele Regresse durchsetzen zu m\"ussen.

Das Pr\"ufsystem ist nicht nur ein schlechtes Messinstrument --- es ist ein hochwirksames Disziplinierungsinstrument, das gerade \textit{weil} es nicht pr\"azise misst, seine gr\"o\ss{}te Wirkung entfaltet.

\textbf{Evidenzbasis:} Theoriegest\"utzte Analyse (Foucault, \textit{Surveiller et punir}, 1975), angewandt auf die Ribbat-Befunde und die Verfahrensstrukturanalyse.

\textbf{Einschr\"ankung:} Theoretische Einordnung, nicht empirischer Nachweis der Intentionalit\"at.

%% ============================================================
\section{Vorgeschlagene Ma\ss{}nahmen (KNV-Tabelle)}
%% ============================================================

Die Studie bewertet 14 Ma\ss{}nahmen nach Aufwand, politischer Realisierbarkeit, erwartetem Nutzen und Kosten-Nutzen-Verh\"altnis (KNV). Die Rangfolge:

\smallskip
{\small
\begin{tabularx}{\textwidth}{c c X c c c c c}
\hline
\textbf{Rg} & \textbf{\#} & \textbf{Ma\ss{}nahme} & \textbf{Aufw.} & \textbf{Real.} & \textbf{Nutz.} & \textbf{KNV} & \textbf{Zeit} \\
\hline
1 & 1 & \textbf{Halbsatz-Streich \S\,106b Abs.\,2 S.\,4} (Beratungsschutz f\"ur V2) & ger. & hoch & hoch & \textbf{A+} & kurz \\
2 & 2 & Sichtbarmachung \S\,48 BMV-\"A im SGB\,V & ger. & hoch & mittel & \textbf{A} & kurz \\
3 & 3 & Bundeseinheitliche Bagatellgrenze & ger. & hoch & mittel & \textbf{A} & kurz \\
4 & 4 & Positive R\"uckkopplung (Unauff.-Zertifikat) & ger. & hoch & mittel & \textbf{A} & kurz \\
5 & 5 & Unabh. Effizienzpr\"ufung des Pr\"ufsystems & ger. & hoch & hoch & \textbf{A} & mittel \\
6 & 6 & Publikationspflicht der Pr\"ufungsstellen & ger. & mittel & transf. & \textbf{A$-$} & kurz \\
7 & 7 & \textbf{PP-003: Regress-Transparenzportal} & mittel & mittel & hoch & \textbf{B+} & mittel \\
8 & 8 & \"Offnung der operativen Pr\"ufkataloge & mittel & mittel & hoch & B+ & mittel \\
9 & 9 & Kostenfolgen-Pflichtabsch\"atzung bei Antr\"agen & mittel & mittel & mittel & B & mittel \\
10 & 10 & Maximale Gesamtverfahrensdauer & mittel & mittel & mittel & B & mittel \\
11 & 11 & Open-Source-Verordnungstool & mittel & mittel & mittel & B & mittel \\
12 & 12 & Sanktionsfolge f\"ur unbegr\"undete Kassenantr\"age & mittel & niedr. & hoch & B & lang \\
13 & 13 & IFG-Anfragen an Pr\"ufungsstellen & ger. & niedr. & ger.--m. & C & kurz \\
14 & 14 & Governance-Pharmakovigilanz & hoch & niedr. & transf. & B$-$ & lang \\
15 & 15 & \textbf{Versorgungspr\"ufung (\S\S\,70, 135a SGB\,V)} --- Messinstrument f\"ur Unter-/Fehlversorgung (symmetrische Gegenpr\"ufung zur WiP) & hoch & mittel & transf. & B & lang \\
\hline
\end{tabularx}
}
\smallskip

\textbf{Erl\"auterung der wichtigsten Ma\ss{}nahmen:}

\begin{itemize}
  \item \textbf{Rang 1 (Halbsatz-Streich):} Eine Zeile Gesetzestext w\"urde die strukturelle Asymmetrie zwischen V1 und V2 beseitigen. Damit g\"alte auch f\"ur Einzelfallpr\"ufungen, dass beim erstmaligen Anlass eine Beratung statt eines Regresses erfolgt. \textbf{Dogmatisch wichtig:} Die Beratung nach \S\,106b Abs.\,2 S.\,2 SGB\,V ist eine \textit{Regelungsentscheidung mit Verwaltungsaktqualit\"at} (nicht nur ein informeller Hinweis; BSG B\,6\,KA\,21/15\,R, B\,6\,KA\,1/17\,R). Sie begr\"undet einen Vertrauenstatbestand und schlie\ss{}t sp\"atere Regresse f\"ur denselben Sachverhalt aus. Bereits 2015 in der KBV-Kritik vorgeschlagen.
  \item \textbf{Rang 2 (\S\,48 BMV-\"A):} Der wichtigste existenzvernichtende Regressgrund (rund 490.000\,EUR Stempelregress, BSG B\,6\,KA\,9/24\,R vom 27.08.2025) lebt au\ss{}erhalb des SGB\,V in einem Kollektivvertrag. Wer \glqq{}nur\grqq{} das Gesetz liest, \"ubersieht ihn. Rein redaktionelle \"Anderung, kein Akteur verliert.
  \item \textbf{Rang 5 (Effizienzpr\"ufung):} Kein Akteur (BRH, BAS, SVR, Bundestag) pr\"uft die Wirtschaftlichkeit der Wirtschaftlichkeitspr\"ufung. Ein gesetzlicher Evaluationsauftrag alle 5 Jahre w\"urde die institutionelle Selbsterhaltungsdynamik durchbrechen.
  \item \textbf{Rang 6 (Publikationspflicht):} Ohne ver\"offentlichte Vollzugsstatistiken ist jede Effizienzdiskussion unm\"oglich. Die spieltheoretische Analyse identifiziert den Abbau asymmetrischer Information als wirkungsm\"achtigste Einzelintervention.
\end{itemize}

%% ============================================================
\section{Bedeutung f\"ur das Regress-Transparenzportal (PP-003)}
%% ============================================================

PP-003 positioniert sich im oberen Mittelfeld der Ma\ss{}nahmenrangliste (\textbf{Rang 7 von 14}, KNV: B+).

\textbf{Strategische Funktion:} PP-003 ist der \glqq{}Plan B, der Plan A erzwingen soll.\grqq{} Wird die Publikationspflicht der Pr\"ufungsstellen (Rang 6) politisch realisiert, wird das Portal teilweise \"uberfl\"ussig --- und genau das w\"are das gew\"unschte Ergebnis. Das Portal erzeugt den informationellen Druck, die Reformhebel den institutionellen. Eines ohne das andere verpufft.

\textbf{Was PP-003 leisten kann:}
\begin{itemize}
  \item Erstmals verfahrensgetrennte (V1 vs. V2) Daten aus \"arztlicher Perspektive erheben
  \item Die Risikokalibrierung der \"Arzte korrigieren (Aggregat-Zahlen statt Anekdoten)
  \item Politischen Reformdruck aufbauen durch Sichtbarmachung der Datenl\"ucke
  \item Zivilgesellschaftlichen Druckaufbau f\"ur die Ma\ss{}nahmen der R\"ange 1--6 leisten
\end{itemize}

\textbf{Was PP-003 nicht leisten kann:}
\begin{itemize}
  \item Die strukturelle Asymmetrie zwischen V1 und V2 beseitigen (daf\"ur braucht es Rang 1)
  \item Die fehlende Effizienzaufsicht ersetzen (daf\"ur braucht es Rang 5)
  \item Das Compliance-Theater auf Kassenseite verhindern (Transparency Paradox nach Bernstein 2012)
  \item Den kausalen Schluss Pr\"ufdruck $\to$ Unterversorgung empirisch schlie\ss{}en
\end{itemize}

\textbf{Robustheit gegen\"uber Kritik:} PP-003 ber\"ucksichtigt das Transparency Paradox (Bernstein 2012) durch drei Gestaltungsentscheidungen: (1)~Wahlkabinen-Prinzip (kryptografische Trennung), (2)~$k$-Anonymit\"at ($k\geq5$), (3)~Fokus auf Aggregate statt Einzelf\"alle. Dennoch ist das Portal ein notwendiger, aber nicht hinreichender Schritt.

%% ============================================================
\section{Was am ehesten umgesetzt werden sollte}
%% ============================================================

\subsection*{H\"ochste Priorit\"at (R\"ange 1--6): Geringe Aufw\"ande, hohe Realisierbarkeit, kurzfristig umsetzbar}

Die R\"ange 1--6 teilen drei Eigenschaften: Sie kosten wenig, sind politisch anschlussf\"ahig und k\"onnen innerhalb eines Jahres wirksam werden.

\textbf{Der effizienteste Einzelhebel} ist der Halbsatz-Streich (Rang 1): Eine Zeile Gesetzestext w\"urde den wichtigsten einzelnen Treiber der Regressangst eliminieren --- die Tatsache, dass bei der Einzelfallpr\"ufung im Gegensatz zur Auff\"alligkeitspr\"ufung kein \glqq{}Beratung vor Regress\grqq{} vorgesehen ist.

\textbf{Der transformativste Hebel} ist die Publikationspflicht der Pr\"ufungsstellen (Rang 6): Solange keine Vollzugsstatistiken existieren, ist jede Reformdebatte eine Debatte ohne Datenbasis.

\textbf{Die politisch einfachsten Hebel} sind die R\"ange 2--4: rein redaktionelle \"Anderungen (Sichtbarmachung \S\,48 BMV-\"A), bereits in KBV-Forderungslisten enthaltene Ma\ss{}nahmen (Bagatellgrenze) oder KV-interne Ma\ss{}nahmen ohne Gesetzes\"anderung (Unauff\"alligkeits-Zertifikat).

\subsection*{Mittlere Priorit\"at (R\"ange 7--11): Erg\"anzende Ma\ss{}nahmen mit mittlerem Aufwand}

PP-003 (Rang 7) entfaltet seine volle Wirkung erst im Verbund mit den R\"angen 1--6. Als zivilgesellschaftlicher Druckaufbau ist es besonders dann wertvoll, wenn die politischen Hebel nicht zeitnah gezogen werden.

\subsection*{Langfristige Vision (R\"ange 12--14): Strukturelle Reformen}

Die Governance-Pharmakovigilanz (Rang 14) --- eine regelm\"a\ss{}ige, unabh\"angige Bewertung der Versorgungs-Nebenwirkungen regulatorischer Ma\ss{}nahmen --- w\"are die umfassendste System\"anderung. Sie hat den h\"ochsten transformativen Nutzen, aber die niedrigste politische Realisierbarkeit.

%% ============================================================
\section{Offengebliebene Fragestellungen}
%% ============================================================

\subsection*{6.1 Fehlende V2-Pipeline-Daten}

F\"ur die Einzelfallpr\"ufung (V2) existiert keine publizierte Pipeline-Statistik. Weder die Gesamtzahl der Antr\"age, noch die Verwerfungsquote, noch die Vergleichsquote, noch die Verfahrensdauer wird \"offentlich publiziert. Die gesamte Spalte \glqq{}k.\,D.\grqq{} (keine Daten) in den Pipeline-Tabellen der Studie wartet auf F\"ullung.

\subsection*{6.2 IFG-Anfragen (Frist ca. 10.05.2026)}

Neun IFG-Anfragen an BMG, BAS, KBV, GPE und weitere Stellen wurden am 10.04.2026 eingereicht. Die Antwortfrist betr\"agt 30 Tage (\S\,7 Abs.\,5 IFG). Erste Empfangsbest\"atigungen liegen vor. Die wichtigsten erwarteten Erkenntnisse:

\begin{itemize}
  \item \textbf{A1 + A2 (BMG, GPE RLP):} W\"urden erstmals die V2-Pipeline mit realen Zahlen f\"ullen --- der \glqq{}gr\"o\ss{}te einzelne Qualit\"atssprung\grqq{} f\"ur die Studie.
  \item \textbf{A3 (BAS):} W\"urden den Niskanen-Befund (GPE-Selbsterhaltung) empirisch untermauern oder widerlegen.
  \item \textbf{A6 (BRH):} Die Antwort \glqq{}keine Sonderpr\"ufung seit 2005\grqq{} w\"are selbst ein Befund.
  \item \textbf{A9 (BMG):} W\"urde erzwingen, dass das BMG schriftlich zu konkreten Schwachstellen im Normgef\"uge Stellung nimmt.
\end{itemize}

Auch Verweigerungen sind publizistisch verwertbar: \glqq{}Diese Daten liegen uns nicht vor\grqq{} best\"atigt die diagnostizierte Transparenzl\"ucke.

\subsection*{6.3 Kausalit\"atsl\"ucke (Pr\"ufdruck $\to$ Unterversorgung)}

Der kausale Schluss \glqq{}Pr\"ufdruck $\to$ ver\"andertes Verordnungsverhalten $\to$ Unterversorgung $\to$ Patientensch\"aden\grqq{} ist nicht empirisch geschlossen. Die Studie verf\"ugt nicht \"uber eine kausalanalytische Studie, die Versorgungsergebnisse direkt mit Pr\"ufintensit\"at korreliert --- etwa einen Vergleich von KV-Regionen mit unterschiedlicher Einzelfallpr\"ufungsfrequenz. Ein solches nat\"urliches Experiment w\"are der st\"arkste Beleg, existiert aber nicht.

\subsection*{6.4 Ribbat-Replikation ausstehend}

Ribbat et al. (2023) ist die einzige gro\ss{}e deutsche Befragungsstudie, die spezifisch die Auswirkungen der Regressangst auf das Verordnungsverhalten misst. Eine unabh\"angige Replikation --- idealerweise mit verfahrensgetrennter Erhebung (V1 vs. V2) --- steht aus. Zwei empirische Befunde w\"urden die Kausalkette erheblich schw\"achen:

\begin{enumerate}
  \item Eine Replikation mit weniger als 20\,\% berichtetem Verhaltenseffekt
  \item Eine Korrelationsanalyse zwischen Pr\"ufintensit\"at und Verordnungsvolumen, die keinen Zusammenhang zeigt
\end{enumerate}

\subsection*{6.5 Folgekostenrechnung empirisch nicht fundiert}

Die Hochrechnung (0,9--1,7 Mrd. EUR/Jahr) beruht auf Szenario-Annahmen. Insbesondere die Attributionsrate (5\,\%) und die Frequenz unterlassener Verordnungen sind Modellannahmen. Die Modellrechnung ersetzt keine kausalanalytische Studie. Eine systematische Pr\"ufung durch eine zust\"andige Stelle (BRH, SVR, IGES) hat bisher nicht stattgefunden.

\subsection*{6.6 Abschreckungswirkung nicht quantifiziert}

Die Studie rechnet nur Kosten (Unterversorgung, Defensivmedizin, Folgesch\"aden), nicht den Nutzen des Pr\"ufsystems (verhinderte Mehrausgaben durch Abschreckung). Diese Abschreckungswirkung ist real und schwer messbar. Eine vollst\"andige Kosten-Nutzen-Analyse m\"usste sie im Nenner ber\"ucksichtigen.

\subsection*{6.7 V2a/V2b-Verteilung unbekannt}

Welcher Anteil der Einzelfallpr\"ufungen V2a (unwirtschaftlich) und welcher V2b (unzul\"assig/Formfehler) betrifft, ist nicht \"offentlich bekannt. Ohne diese Differenzierung kann die Sch\"arfe des Befunds nicht kalibriert werden.

\subsection*{6.8 Automatisiertes Screening durch Kassen --- offene Rechtsfrage}

Ob Krankenkassen tats\"achlich systematisch alle Verordnungsdaten mit Pr\"ufsoftware scannen, um daraus Einzelfallpr\"ufungsantr\"age zu generieren, ist empirisch nicht belegt. Belegt ist, dass die technischen Mittel existieren und dass das Volumen (KV Nordrhein: ca. 20.000 Antr\"age/Jahr) schwer allein mit zuf\"alligen Verdachtsmomenten erkl\"arbar ist. Die verfassungsrechtliche Einordnung (BVerfG Hessendata 2023, 1\,BvR\,1547/19) ist argumentativ vertretbar, aber gerichtlich nicht gekl\"art.

\subsection*{6.9 GPE-Kosten extern nicht verifizierbar}

Die Kosten der bundesweiten Pr\"ufinfrastruktur (ca. 600--900 Stellen) sind aus offenen Quellen nicht verifizierbar. Die GPE sind keine staatlichen Beh\"orden, tauchen in keinem Haushalt auf, und der Bundesrechnungshof pr\"uft sie nicht. Diese strukturelle Nicht-Verifizierbarkeit ist selbst ein Befund.

%% ============================================================
\section{Methodische Transparenz}
%% ============================================================

Die Studie entstand in einem mehrstufigen, KI-gest\"utzten Analyseprozess im April 2026 (gesch\"atzte Gesamtarbeitszeit ca. 60--80 Stunden, Mensch + KI kombiniert). 31 eigenst\"andige Recherchestr\"ange wurden durchgef\"uhrt und in eine integrierte Argumentation \"uberf\"uhrt.

\textbf{Multi-Modell-Architektur:} Vier Instanzen arbeitsteilig --- Claude (Opus 4.6, Prim\"aranalyse und Synthese), Copilot (GPT-4o, Konzeptsch\"arfung und Fehlkonzept-Erkennung), Gemini (Deep Research, Theorierahmen und Quellenerg\"anzung), LG/Mensch (Fragestellungen, Richtungsentscheidungen, finale Freigabe).

\textbf{Bias-Management:} Cross-Model-Divergenzen wurden systematisch als Signal genutzt. F\"unf dokumentierte Hypothesenrevisionen (u.\,a. Datierung der Einzelfallpr\"ufung: GSG 1992 $\to$ GRG 1988; Moosazadeh-Fehlzuordnung $\to$ Zheng 2023; Kostenmodell: Systemrahmen $\to$ Dreistufenmodell nach Copilot-Review). Ein eigener Devil's-Advocate-Agent versuchte systematisch, die zentrale These zu widerlegen.

\textbf{Cross-Source-Divergenz (Phase XIV):} 15 spezialisierte Agenten (6 Opus, 9 Sonnet) recherchierten parallel mit 25 Datenpunkten aus 14 unabh\"angigen Domains. Zentrale Lektion: Cross-Model-Divergenz kontrolliert Modell-Bias, aber nicht Quellen-Bias. Erst Multi-Quellen-Recherche brachte die SpiFa-Bayern- und BMG-Daten ans Licht. Zwei Prior-Art-Checks best\"atigten: Beide zentralen Synthesen (Definitionsdivergenz; Evidenz in formaler Verkleidung) sind origin\"ar.

\textbf{Halluzinationskontrollen:} Gemini-Outputs erforderten systematischen Double-Check bei aktuellen politischen Ereignissen (dokumentierte F\"alle: Moosazadeh-Fehlzuordnung, franz\"osische Protestzahlen). Jede Quelle wurde gegen die Originalpublikation gepr\"uft.

\textbf{Grenzen der Evidenz:} Die Studie behauptet \textbf{nicht}, dass das Wirtschaftlichkeitspr\"ufsystem abgeschafft werden sollte, dass alle Regresse ungerechtfertigt sind, oder dass \"Arzte keiner Kontrolle bed\"urfen. Sie behauptet:

\begin{enumerate}
  \item Das System misst \"uberwiegend nicht, was es zu messen vorgibt (Zertifikatsfehler) --- wobei ${\sim}$35\,\% der V2-Pr\"ufungen durchaus evidenzbasierte Steuerung leisten. Das Problem ist nicht die Evidenz, sondern der Kommunikationskanal: punitiv statt edukativ.
  \item Die Einzelfallpr\"ufung ist strukturell asymmetrisch ausgestaltet (kein Beratungsschutz, normativer Schadensbegriff bei V2b).
  \item Die ``unter 1\,\%''-Zahl und die 72\,\%-Erfahrungsrate messen verschiedene Dinge (Definitionsdivergenz, Faktor ${\sim}$500$\times$). Die Diskrepanz ist ein Messproblem, kein Wahrnehmungsfehler.
  \item Die Datenlage ist so d\"unn, dass weder Bef\"urworter noch Kritiker ihre Position empirisch vollst\"andig belegen k\"onnen.
  \item Die Transparenzl\"ucke ist nicht zuf\"allig, sondern strukturell --- und sie zu schlie\ss{}en w\"are der wirksamste einzelne Schritt zur Versachlichung der Debatte.
\end{enumerate}

Wo die Studie Szenario-Annahmen macht, sind diese als solche gekennzeichnet. Wo Wertungen vorgenommen werden, sind sie als Wertungen kenntlich. Wo Befunde mehrere Lesarten zulassen, werden alle Lesarten dargestellt.

\vspace{4mm}
\noindent\rule{\textwidth}{0.4pt}

\vspace{2mm}
\noindent\textbf{Vollst\"andige Studie:} ST-001 Systemtheoretische Aufarbeitung der Regressangst (Arbeitsversion v0.14, April 2026)\\
\textbf{Begleitstudie zu:} PP-003 Regress-Transparenzportal (v2.9)\\
\textbf{Lizenz:} CC BY 4.0\\
\textbf{Kontakt:} \href{mailto:hallo@um-bruch.org}{hallo@um-bruch.org} \quad|\quad um-bruch.org\\
\textbf{Repository (Quellen, Recherche, Versionen):} \href{https://github.com/research-line/regressangst}{github.com/research-line/regressangst}\\
\textbf{Companion-Software (Open Source):} \href{https://github.com/research-line/verordnungsampel}{github.com/research-line/verordnungsampel} \textit{-- VerordnungsAmpel, Werkzeug zur Regress-Risikoindikatoren-Anzeige ohne Gew\"ahr.}

\vspace{4mm}
\noindent\rule{\textwidth}{0.4pt}

\vspace{2mm}
\noindent\textbf{Änderungshistorie:}\\
\textbf{v1.0} (April 2026) --- Erstversion, basierend auf ST-001 v0.11. 15 Befunde, KNV-Tabelle (14 Maßnahmen), PP-003-Einordnung, 9 offene Forschungsfragen.\\
\textbf{v1.0+CORE4} (April 2026) --- CORE4-Integration: 2 neue Befunde (Definitionsdivergenz mit Drei-Schichten-Tabelle; Doppelfunktion mit A/B/C-Kategorien). V2-Datenabschnitt korrigiert (BMG 47.000, SpiFa Bayern 13.332). ``Grenzen der Evidenz'' auf 5 Punkte erweitert. Multi-Agenten-Experiment + Prior-Art-Checks in Methodische Transparenz. ST-001-Referenz auf v0.12 (87~Seiten). 13~Seiten.\\
\textbf{v1.1+CORE4} (April 2026) --- Synchronisation mit ST-001 v0.14 (95~Seiten) nach 18-facher Review-Chain: Poisson-Bandbreite $\lambda = 0{,}30$--$1{,}48$ mit 60--99\,\%-Gesamtbandbreite erg\"anzt; Beratung als Regelungsentscheidung (\S\,106b Abs.\,2 S.\,2 SGB\,V; BSG B\,6\,KA\,21/15\,R, B\,6\,KA\,1/17\,R) pr\"azisiert; \"ubergeordneter Zertifikatsfehler und SVR-Trias (\S\S\,70, 135a SGB\,V) verankert; Versorgungspr\"ufung als neue KNV-Zeile~15; Folgekosten-Detaillierung auf mittleres Szenario $\approx 1{,}1$\,Mrd.\,EUR; BSG B\,6\,KA\,9/24\,R Datum (27.08.2025) und Summe (rund 490.000\,EUR) vereinheitlicht.\\
\textbf{v1.2} (April 2026, 2026-04-14) --- Synchronisation mit ST-001 v0.15 (P1-Block): \textbf{Genre-Kl\"arung} in der Zusammenfassung --- ``multiperspektivische Bestandsaufnahme mit Pilot-Charakter'' statt Original-Studie; \textbf{quantitative Einsch\"atzung} in Abschnitt V2a/V2b: 60--65\,\% der V2-Verfahren (Kategorie~B + verfahrensfremd) unter zweitem Zertifikatsfehler, 35--40\,\% (Kategorie~A) codieren G-BA-Evidenz punitiv statt edukativ (Devil's-Advocate-Befund aus Einwand~5 in Hauptthese). Untertitel und ST-001-Referenz auf v0.15 angehoben.

\vspace{2mm}
\noindent\textbf{KI-Transparenzhinweis:} Dieses Dokument wurde unter extensivem Einsatz von KI-Sprachmodellen (u.\,a. Anthropic Claude, Microsoft Copilot, Google Gemini) erarbeitet. Es wurden mehrere Kontroll- und Reviewzyklen durchgef\"uhrt~--- KI-gest\"utzt und menschlich. KI-Modelle neigen zu Halluzinationen; Fehler sind trotz sorgf\"altiger Pr\"ufung nicht ausgeschlossen. Wir korrigieren dokumentierte Fehler und ver\"offentlichen verbesserte Versionen. Dieses Dokument stellt keine Rechtsberatung dar. Hinweise: \href{mailto:hallo@um-bruch.org}{hallo@um-bruch.org}

\end{document}
